This is a small CPU study and its scope is narrow. (i) One dataset (8$\times$8 digits), one architecture (a two-hidden-layer MLP) and two synthetic shift families; rotation of $8{\times}8$ images is destructive, so accuracy drops to chance at r45 and r60, and the conclusions may not transfer to natural corruptions, CNNs or larger images. (ii) Five seeds (three for the sweeps), a 450-image test set per seed and paired $t$-tests without multiple-comparison correction; p-values are descriptive, and the H1 non-result may be a lack of power. (iii) Baselines are reimplemented, not tuned: no weight decay, early stopping or per-method hyperparameter search, and the MC-dropout network differs from the other systems in training regularisation, which confounds the system comparison (Section~\ref{sec:ablations}); the dropout decomposition was exploratory and added after seeing results. (iv) ECE is a binned estimator that is easily dominated by the mean confidence when accuracy is near chance; NLL and Brier are shown alongside but the three metrics are strongly correlated across conditions here. (v) The oracle temperature uses the shift itself and 269 validation images, and achieves low ECE with nearly uniform predictions; it shows a mechanism, not a practical method. (vi) Models are cached per seed so rows across shifts are not independent trainings; shift noise is a single draw per seed. A fuller study would add CNNs and corrupted natural-image benchmarks, regularised and larger ensembles, post-hoc methods designed for shift, more seeds, and a proper OOD-detection evaluation.
